\chapter{Sharding}
\label{chap:sharding}

\section{Shard-per-core direction}

\texttt{ShardedExecutor} connects the single-threaded executor to the
shard-per-core model. It starts one executor per shard thread and provides APIs
for submitting work to one shard, all shards, or a map/reduce over shards.

Work does not migrate implicitly. A future runs on its assigned shard until
code deliberately submits work elsewhere.

Parallelism comes from independent shards rather than threads mutating one
service object. Each shard owns its executor, queues, timers, readiness
registrations, and local values. Cross-shard communication is therefore an
explicit cost and ownership boundary.

\begin{center}
\begin{tikzpicture}[
    shardbox/.style={draw=sitasMuted, fill=sitasCodeBackground,
        minimum width=28mm, minimum height=7mm, font=\footnotesize\ttfamily},
    shardlabel/.style={font=\footnotesize\bfseries}
]
\foreach \x/\name in {0/Shard 0, 1/Shard 1, 2/Shard 2} {
    \begin{scope}[xshift=\x*35mm]
        \node[shardlabel] at (0, 2.0) {\name{} thread};
        \node[shardbox] at (0, 1.3) {Executor};
        \node[shardbox] at (0, 0.5) {local tasks};
        \node[shardbox] at (0, -0.3) {local timers};
        \node[shardbox] at (0, -1.1) {local state};
    \end{scope}
}
\end{tikzpicture}
\end{center}

\section{Cross-shard submission}

\texttt{ShardedSubmitter} is the explicit cross-shard async mechanism. A task
on one shard can submit a future to another shard and await the returned join
handle. The remote future is polled on the target shard. The awaiting task
resumes on its original shard after the remote work completes.

Submitters own spawner clones. They are lifetime capabilities and must be
dropped before a runtime can fully drain.

If shard 0 submits work to shard 1, shard 1 polls the submitted future. Its
owned result completes the join handle, waking the awaiting task on shard 0;
the awaiting task does not migrate.

\begin{center}
\begin{tikzpicture}[
    sbox/.style={draw=sitasMuted, fill=sitasCodeBackground,
        minimum width=55mm, minimum height=8mm, font=\footnotesize\ttfamily,
        rounded corners=1.5pt},
    arr/.style={->, >=stealth, thick, color=sitasAccent},
    lbl/.style={font=\tiny\itshape, color=sitasMuted}
]
\node[sbox] (s0a) at (-1, 0) {Shard 0: submit future};
\node[sbox] (s1)  at (6, 0) {Shard 1: poll future};
\node[sbox] (s1r) at (6, -1.4) {Shard 1: owned result};
\node[sbox] (s0b) at (-1, -1.4) {Shard 0: receive result};

\draw[arr] (s0a) -- node[above, lbl] {Send + 'static} (s1);
\draw[arr] (s1) -- (s1r);
\draw[arr] (s1r) -- node[above, lbl] {owned value} (s0b);
\end{tikzpicture}
\end{center}

This is why cross-shard APIs require \texttt{Send + 'static} where a future or
value can move to another thread. It is also why references into shard-local
state cannot be submitted.

The submitter does not steal work. Submitting to shard 1 explicitly states that
shard 1 is the correct execution context.

\section{Mailbox routing}

\texttt{ShardedSubmitter} moves futures. The host-only
\texttt{shard\_mailbox} module moves owned messages through bounded queues. A
mailbox transfers ownership to one shard-local receiver; it does not make the
receiver's application state shared. The portable \texttt{ShardRuntime}
channel has the same cloneable-sender, single-receiver shape, but uses the
\texttt{no\_std} ring buffer and serializes concurrent producers at push.

There are two routing modes.

\begin{description}
    \item[Uniform shard routing] A key maps directly to a physical
          \texttt{ShardId}. The target is one inbound mailbox per shard:
          \texttt{ShardMailboxSet<M>}.
    \item[Non-uniform work-unit routing] A key maps to a logical work-unit
          name. A placement table maps that name to the shard that owns its
          receiver: \texttt{WorkUnitMailboxSet<N, M>}. Multiple work units can
          be assigned to the same shard, and some shards can have no work unit
          for a particular protocol.
\end{description}

\begin{center}
\begin{tikzpicture}[
    routebox/.style={draw=sitasMuted, fill=sitasCodeBackground,
        minimum width=34mm, minimum height=8mm, text width=32mm,
        align=center, font=\scriptsize\ttfamily, rounded corners=1.5pt},
    label/.style={font=\footnotesize\bfseries, color=sitasMuted},
    arr/.style={->, >=stealth, thick, color=sitasAccent}
]
\node[label] (uniform-label) at (-4, 0.7) {Uniform};
\node[routebox] (uniform-key) at (-2.6, 0) {key};
\node[routebox] (uniform-shard) at (1.2, 0) {ShardId};
\node[routebox] (uniform-set) at (5.0, 0) {ShardMailboxSet\\<M>};

\node[label] (nonuniform-label) at (-3.6, -0.9) {Non-uniform};
\node[routebox] (nonuniform-key) at (-2.6, -1.6) {key};
\node[routebox] (work-unit) at (1.2, -1.6) {WorkUnitName};
\node[routebox] (work-set) at (5.0, -1.6) {WorkUnitMailboxSet\\<N, M>};
\node[routebox] (assigned) at (8.8, -1.6) {assigned ShardId};

\draw[arr] (uniform-key) -- (uniform-shard);
\draw[arr] (uniform-shard) -- (uniform-set);

\draw[arr] (nonuniform-key) -- (work-unit);
\draw[arr] (work-unit) -- (work-set);
\draw[arr] (work-set) -- (assigned);
\end{tikzpicture}
\end{center}

\texttt{UniformShardRouter} makes the first path explicit. It routes a
hashable key to a \texttt{ShardId}. \texttt{WorkUnitRouter<N>} makes the
second path explicit. It routes a hashable key to a logical name. In both
cases the key is still the discriminator; what changes is the target type.

The receiving side remains shard-local. A named work unit has exactly one
receiver, and \texttt{receiver\_for\_current\_shard} checks that the receiver is
taken by a task on the work unit's assigned shard. For non-uniform placement,
the logical name alone is not the owner; the name plus its placement determines
which executor may mutate the associated state.

\begin{notabene}[Mailbox routing is not load balancing]
The mailbox layer does not steal work and does not rebalance a hot work unit, but
it makes routing explicit. Changing the cluster shape means changing the
router targets or the work-unit placement table, then observing the new
behavior through owned snapshots and example output.
\end{notabene}

\section{Cache behavior of shard movement}

Shared-nothing ownership can improve cache locality, but messages and queue
metadata still cause hardware traffic.

A shard owns its mutable application state. If a key-value shard repeatedly
updates its local map, the working set for that map tends to stay in the cache
hierarchy of the core running that shard. Other shards do not acquire a mutex
around the map, read its buckets, or update its counters. This avoids a common
source of coherence traffic: several cores repeatedly invalidating the same
cache lines while trying to mutate one logical service object.

Cross-shard communication is different. Sending a message transfers ownership
at the Rust and runtime boundary, but the bytes still move through the machine.
Suppose shard 0 builds an owned message and sends it to shard 1:

\begin{enumerate}
    \item Shard 0 writes the message payload and queue slot.
    \item Shard 0 updates producer-side queue metadata.
    \item Shard 1 observes that the queue is no longer empty.
    \item Shard 1 reads the message payload and mutates its own local state.
\end{enumerate}

The payload is likely warm on shard 0 after construction. Shard 1 must fetch
those cache lines before consuming it. Rust ownership prevents stale borrowed
access on shard 0; it does not remove the physical transfer.

Enqueue and dequeue also touch coordination state such as queue positions,
capacity, close flags, wakers, and counters. These fields are transport
synchronization rather than shared application state, but they can still
become cache-coherence hot spots.

\begin{center}
\begin{tikzpicture}[
    box/.style={draw=sitasMuted, fill=sitasCodeBackground,
        minimum width=40mm, minimum height=8mm, font=\footnotesize\ttfamily,
        rounded corners=1.5pt},
    metabox/.style={draw=sitasAccent, fill=sitasCodeBackground,
        minimum width=40mm, minimum height=8mm, font=\footnotesize\ttfamily,
        rounded corners=1.5pt},
    arr/.style={->, >=stealth, thick, color=sitasAccent},
    lbl/.style={font=\tiny\itshape, color=sitasMuted}
]
\node[box] (s0) at (-5.5, 0) {Shard 0 cache};
\node[metabox] (q) at (0, 0) {mailbox metadata};
\node[box] (s1) at (5.5, 0) {Shard 1 cache};
\node[box] (payload0) at (-5.5, -1.4) {owned payload};
\node[box] (payload1) at (5.5, -1.4) {owned payload};

\draw[arr] (s0) -- node[above, lbl] {enqueue / wake} (q);
\draw[arr] (q) -- node[above, lbl] {dequeue} (s1);
\draw[arr] (payload0) -- node[below, lbl] {cache-line transfer} (payload1);
\end{tikzpicture}
\end{center}

This is the distinction the architecture relies on:

\begin{description}
    \item[Application state] remains local and is mutated only by the owning
          shard. This is where the model avoids uncontrolled shared-cache-line
          contention.
    \item[Transport state] is allowed to synchronize. Mailboxes, reply handles,
          wakers, lifecycle flags, and counters may use atomics or locks
          internally. Their cache behavior matters, but it is confined to the
          runtime boundary.
    \item[Transferred values] are owned and type-safe, but they still occupy
          cache lines and memory bandwidth when consumed on another shard.
\end{description}

\subsection{Payload size and batching}

Small commands can create high metadata traffic at high rates; large messages
can evict useful receiver data. The balance is workload-dependent.

For hot paths, prefer commands that describe the operation compactly and let
the destination shard perform the mutation locally. When many tiny records
share the same route, batching them into a \texttt{Vec<T>} can amortize queue
metadata updates and wakeups. Batching should still preserve ownership: the
batch is an owned value, and after receipt only the destination shard mutates
the state associated with it.

Large buffers need a different tradeoff. Sending an owned buffer is semantically
clean, but the receiving shard will pull the buffer's cache lines into its own
cache hierarchy. If the producing shard also needs the same data again, the
program may be better served by routing the work to the shard that already owns
or will next use the data, rather than bouncing a large working set across
cores.

\subsection{False sharing}

False sharing occurs when independent fields used by different cores reside on
the same cache line. For example, a producer-updated counter and a
consumer-updated counter may be logically separate but physically adjacent. If
they share a cache line, updates by one core can invalidate the other core's
copy of the line.

This can affect:

\begin{itemize}
    \item mailbox producer and consumer indices;
    \item per-shard progress counters;
    \item wake flags and closed flags;
    \item high-frequency observability counters;
    \item arrays of per-shard statistics where adjacent shards update adjacent
          entries.
\end{itemize}

Padding or cache-line alignment is sometimes justified for fields that are
known to be hot and written by different cores. It should be applied
selectively: padding every structure makes memory use worse and can reduce
cache density for data that is mostly read or only touched by one shard.

\subsection{Many producers and one receiver}

An inbound mailbox per shard is clear and predictable, but it can still become
a coherence hotspot. If many shards send to one target at high frequency, the
target's inbound queue metadata becomes shared by many producers. The
application state behind the receiver is still protected by the ownership
model, but the transport path may limit throughput.

Mitigations include:

\begin{itemize}
    \item route data so work is more evenly partitioned;
    \item combine small messages before sending;
    \item use logical work units to split one protocol across multiple
          receivers when the state model allows it;
    \item avoid request patterns where every shard repeatedly asks one shard to
          perform tiny mutations;
    \item return owned summaries rather than repeatedly querying remote state.
\end{itemize}

\subsection{NUMA and placement}

On multi-socket systems, cache and memory movement may cross a NUMA boundary.
The cost of sending a message between two cores on the same socket can differ
from sending between sockets. CPU placement does not by itself solve NUMA
placement, but it makes the question visible: which shard runs where, and what
data does it mostly touch?

CPU and memory placement are explicit Linux requests. Applied CPU placement
records the NUMA node when sysfs exposes it. An optional shard-thread default
memory policy can bind, prefer, or interleave nodes, or select the node local to
the pinned CPU. Required placement converts failure into startup failure;
otherwise snapshots record the result. Cross-shard and especially cross-socket
movement still has nonzero cost after successful placement.

This is not complete allocator control. The policy affects future allocations
made by the shard thread according to kernel rules. The runtime does not yet
expose \texttt{mbind} for existing address ranges, page migration, or
allocator-specific arenas.

\subsection{What the model buys}

The model does not promise zero-copy communication between arbitrary shards and
does not make a mailbox free. Its promise is narrower and stronger:

\begin{itemize}
    \item application state has one mutating owner;
    \item cross-shard movement is explicit in the API;
    \item values crossing the boundary are owned;
    \item runtime synchronization is localized to queues, replies, wakeups, and
          lifecycle bookkeeping;
    \item cache-coherence costs are easier to find because they occur at named
          transfer points rather than hidden shared-state access.
\end{itemize}

Mailboxes are shared transport mechanisms. Application service state still
belongs to the receiving shard.

\section{Shard-local values}

\texttt{ShardLocal<T>} owns one value per executor shard. Access is routed to
the owning shard and the closure receives \texttt{\&mut T} synchronously. The
implementation uses internal unchecked storage plus runtime owner checks, not a
mutex.

References to local state cannot escape the closure and cannot cross
\texttt{.await}. Ownership stays local even when the runtime uses
synchronization for task bookkeeping.

Each shard owns one \texttt{T}; the API must prevent references to that value
from becoming async state.

The API prefers closures returning owned results. If a computation needs to
wait, copy or compute an owned value from the shard-local state first, then
\texttt{await} using that owned value outside the access closure.

\begin{notabene}[Why not async closures for local access?]
An async closure can suspend while holding references captured from its
environment. That is exactly what shard-local state must avoid. The synchronous
closure form is less flexible, but it makes the safe pattern obvious: inspect
or mutate local state now, return an owned value, and only then await.
\end{notabene}

\section{Sharded scheduling groups}

A sharded scheduling group is not a global queue. It is a convenience handle
containing one executor-local scheduling group per shard. Creating a group
called \texttt{foreground} on a four-shard runtime creates four independent
local \texttt{foreground} groups, all with the same shares.

Each shard still schedules only its own tasks. The sharded group handle lets
callers use one typed API to spawn or submit matching classes of work across
shards.

Per-shard snapshots keep the same boundary: each shard's local group poll
count, charged poll time, average poll duration, and poll-time share. A
monitoring tool can compare those owned values across shards, but the runtime
does not introduce a global scheduling queue to produce them.

Runtime identity matters. A sharded group created by one
\texttt{ShardedExecutor} is rejected by another runtime, even if both runtimes
have a group with the same name. This prevents accidental targeting of
same-numbered executor-local group ids in the wrong runtime.

\section{CPU placement}

CPU placement is an experimental Linux-supported runtime request. Linux applies
hard affinity with \texttt{sched\_setaffinity} and observes container cpuset
restrictions through \texttt{sched\_getaffinity}. Non-Linux platforms report
unsupported.

By default, placement failures are recorded in shard snapshots and startup
succeeds. Required placement turns that into fail-fast behavior.

On Linux, the runtime asks the kernel to pin shard threads to particular CPUs.
In containers, the available CPU set may already be restricted, so the runtime
first queries what the kernel allows. On macOS and other platforms, placement
reports unsupported rather than pretending that pinning happened.

Snapshots separate requested placement, the host's allowed CPU set, and the
observed result. When available, they also report the NUMA node associated with
the pinned CPU.

Memory placement follows the same explicit-request model. The default is no
memory policy change. When requested, Linux applies a thread default memory
policy after CPU placement and before the shard executor runs. Snapshot fields
separate CPU placement from memory placement so callers can see whether each
request was applied, unsupported, or rejected.
